\section[Generic proximal geometry]{Generic relative proximal accuracy and canonical vertex ranges}
\label{sec:op3-generic-proximal-geometry}

\paragraph{Scope and source.}
This section concerns a supplied convex objective and supplied Laplacians;
its accuracy parameter is an inner relative-energy tolerance, not
$\epsappr$ or $\rho$. The new statements are \statusproved{} drafts
awaiting independent review. The convergence inequality
\eqref{eq:op3-source-inexact-apg-bound} is imported from CPW Algorithm 9
and Theorem 8.3, PDF p.\ 44 \citep{chen2021diffusion}; the normalization
agrees with Claim 8.19, pp.\ 50--51. We prove the tolerance directly,
without invoking Claim 8.21's all-number-range assumption. This does not
import a fast inner oracle or establish an unrestricted recursive solver.

Let $\Phi(x)=x^\top\mathcal Gx/2+h(x)$, where $h$ is proper closed convex,
$0$ is feasible, a finite minimum $\Phi^*$ is attained, and
$0\preceq\mathcal G\preceq\mathcal H\preceq\kappa\mathcal G$ with
$\kappa\geq1$. Here $\mathcal G,\mathcal H$ may initially be arbitrary
positive semidefinite matrices. Write $\Delta=\Phi(0)-\Phi^*$, and use
the model $\Phi_x$ and normalized objective $E_x$ in
\eqref{eq:op3-capped-prox-model}--\eqref{eq:op3-capped-prox-normalization},
with the general $h$ and with $E_x=\Phi_x-\Phi_x(0)$.

\begin{theorem}[A universal relative proximal tolerance; \statusproved]
\label{thm:op3-generic-relative-proximal}
Choose the least power of two $T$ satisfying $T^2\geq256\kappa$ and set
\begin{equation}
 \delta_{\rm rel}=\frac1{33\cdot2^{30}\kappa^3}.
 \label{eq:op3-generic-relative-policy}
\end{equation}
Run source Algorithm 9 from $y_0=z_0=0$, obtaining a feasible output
$q$ at every center $x$ with
$E_x(q)\leq E_x^*/(1+\delta_{\rm rel})$.
Then every absolute model error is at most
$e_0:=\Delta/(2^{30}\kappa)$ and
\[
 \Phi(y_T)-\Phi^*\leq\Delta/32.
\]
No knowledge of $\Delta$, a positive gap floor, vertex coordinate bounds,
or absolute coefficient scale is needed to set this tolerance.
All iterates obey $\Phi(y_k)-\Phi^*\leq2\kappa\Delta$ and all centers
obey $\|x\|_{\mathcal H}^2\leq64\kappa^2\Delta$.
\end{theorem}
\begin{proof}
Fix an optimum $u$. Convexity and optimality give, for every feasible $v$,
\begin{equation}
 \|v-u\|_{\mathcal H}^2\leq
 2\kappa\bigl(\Phi(v)-\Phi^*\bigr),\qquad
 \|u\|_{\mathcal H}^2\leq2\kappa\Delta.
 \label{eq:op3-generic-gap-distance}
\end{equation}
First suppose $\Delta>0$. Since $T<32\sqrt\kappa$,
$\sum_{j=1}^Tj\leq528\kappa$, and hence
\[
 3\sum_{j=1}^Tj\sqrt{2e_0}
 \leq\frac{1584\sqrt2}{2^{15}}\sqrt{\kappa\Delta}
 <\frac18\sqrt{\kappa\Delta}.
\]
For any already justified prefix of length $k\geq2$, the source bound
therefore gives $\Phi(y_k)-\Phi^*\leq8\kappa\Delta/(k+1)^2$.
At the first step $x=0$, the relative contract gives $E_0(q)\leq0$,
so $\Phi(q)\leq\Phi_0(q)\leq\Phi(0)$. Thus every justified previous
iterate, including $y_0$, has gap at most $2\kappa\Delta$ and distance
at most $2\kappa\sqrt\Delta$ from $u$ in the $\mathcal H$ seminorm.

For completeness the update with $A_{k+1}=(k+2)/2$ is
$x_{k+1}=(1-1/A_{k+1})y_k+z_k/A_{k+1}$ and
$z_{k+1}=z_k+A_{k+1}(y_{k+1}-x_{k+1})$.
It implies $z_k=A_ky_k-(A_k-1)y_{k-1}$ for $k\geq1$, and
\[
 x_{k+1}=y_k+\frac{k-1}{k+2}(y_k-y_{k-1}).
\]
The coefficient lies in $[0,1)$, so
$\|x-u\|_{\mathcal H}\leq6\kappa\sqrt\Delta$ and
$\|x\|_{\mathcal H}\leq8\kappa\sqrt\Delta$; the first center is zero.
Since $\Phi_x\geq\Phi$, this bound uses only earlier errors and gives
\begin{equation}
 0\leq\Phi_x(0)-\Phi_x(q)
 \leq\tfrac12\|x\|_{\mathcal H}^2+\Delta
 \leq33\kappa^2\Delta.
 \label{eq:op3-generic-model-range}
\end{equation}
The relative contract rearranges to
$0\leq e_x\leq\delta_{\rm rel}(\Phi_x(0)-\Phi_x(q))\leq e_0$.
This closes the induction before applying its next prefix bound.
At $T$, that bound is at most $8\kappa\Delta/T^2\leq\Delta/32$.
As in the previous section, the source seminorm inequality can be justified
by applying its positive definite version with $\mathcal H+\nu I$ to
the same fixed finite run and letting $\nu\downarrow0$.

If $\Delta=0$, take $u=0$. Inductively a center in $\ker\mathcal H$
has $\Phi_x(0)=\Phi^*$ and $E_x^*=0$. The relative output is therefore
also an optimum of $\Phi$. Equation~\eqref{eq:op3-generic-gap-distance}
puts it in $\ker\mathcal H$, and the linear updates keep the next center
in that kernel. All gaps and errors are zero, even if the oracle chooses
large coordinates in null directions. This proves the zero-gap case.
\end{proof}

\begin{corollary}[Translation for restarts; \statusproved]
\label{cor:op3-generic-residual-restart}
For a feasible starting vector $a$, apply
\cref{thm:op3-generic-relative-proximal} to
$\Psi(z)=\Phi(a+z)-\Phi(a)$. The same tolerance gives a factor-$32$
gap contraction from $a$. For separable VWFs with domains $[L_i,\infty)$,
the residual functions are
\[
 f_i^a(t)=f_i(a_i+t)-f_i(a_i)+(\mathcal Ga)_it,
 \qquad t\geq L_i-a_i.
\]
They have nonpositive lower endpoints, value zero at zero, continuous
concave derivatives and constant final derivatives. Thus they remain
source VWFs, including when translated breakpoints are negative.
\end{corollary}
\begin{proof}
Expand the quadratic in $\Phi(a+z)$. Translation and addition of an
affine function preserve the stated derivative properties, and
$a_i\geq L_i$. The optimum and initial gap translate exactly.
\end{proof}

We now specialize to a connected supplied graph on $n\geq2$ vertices
with positive edge weights, Laplacian $\mathcal G$, and finite-piece
VWFs $f_i$ on $[L_i,\infty)$, $L_i\leq0$.
Let $T_i$ be each final derivative and put
\begin{equation}
 \begin{aligned}
 U=\max(0,\text{all vertex breakpoints}),\quad
 &B_0=\max(U,\max_i(-L_i)),\\
 A=\sum_i|f_i'(0)|,\quad &K=\frac{n-1}{c_{\min}},
 \end{aligned}
 \label{eq:op3-generic-input-parameters}
\end{equation}
where the maximum over an empty breakpoint list is zero and $c_{\min}$
is the smallest positive supplied edge weight.

\begin{proposition}[Finite minimum and a canonical representative; \statusproved]
\label{prop:op3-vwf-canonical-input-gap}
The objective $\Phi(q)=q^\top\mathcal Gq/2+\sum_i f_i(q_i)$
has a finite infimum if and only if $\sum_iT_i\geq0$. In that case its
minimum is attained. Every feasible $q$ can be replaced by
\begin{equation}
 \mathcal C(q)=q-\max(0,\min_iq_i-U)\mathbf1
 \label{eq:op3-vwf-canonical-map}
\end{equation}
without increasing energy. The new minimum coordinate is at most $U$,
and the explicit initial-gap bound is
\begin{equation}
 0\leq\Phi(0)-\Phi^*\leq B_{\rm gap}:=AB_0+KA^2/2.
 \label{eq:op3-vwf-input-gap-bound}
\end{equation}
If $\mathcal H$ is also a graph Laplacian, the same replacement never
increases any normalized proximal objective $E_x$.
\end{proposition}
\begin{proof}
Along $q=t\mathbf1$ beyond all breakpoints, energy has slope
$\sum_iT_i$; a negative slope gives infimum $-\infty$.
Suppose the slope sum is nonnegative. If the shift in
\eqref{eq:op3-vwf-canonical-map} is positive, all changed coordinates
remain at least $U\geq0\geq L_i$ and stay on their final rays. The
Laplacian energy is unchanged and the vertex sum decreases by the shift
times $\sum_iT_i$. The same is true for $\Phi_x$, because
$(\mathcal H-\mathcal G)\mathbf1=0$.

For a canonical vector write $b=\min_iq_i$ and
$D=\max_iq_i-\min_iq_i$. Then $|b|\leq B_0$. A simple path between
minimum and maximum coordinates and weighted Cauchy--Schwarz give
$q^\top\mathcal Gq\geq D^2/K$. The tangents at zero yield
\begin{align*}
 \Phi(q)&\geq\Phi(0)+D^2/(2K)+\sum_i f_i'(0)q_i\\
 &\geq\Phi(0)-AB_0+D^2/(2K)-AD\\
 &\geq\Phi(0)-AB_0-KA^2/2.
\end{align*}
For a general vector first canonicalize, which only lowers energy, so
the same lower bound holds. A canonical minimizing sequence has bounded
energy, hence bounded $D$ by the penultimate display. Its minimum lies
in the compact interval $[\min_iL_i,U]$, so the whole sequence is bounded.
The feasible domain is closed and the piecewise quadratic objective is
continuous there. A convergent subsequence attains the finite minimum.
This also proves the input-gap bound, including zero total final slope.
\end{proof}

\begin{corollary}[Input-controlled canonical accelerated ranges; \statusproved]
\label{cor:op3-generic-canonical-ranges}
Under the graph hypotheses above, canonicalize every relative-oracle
output before the source updates. The conclusion of
\cref{thm:op3-generic-relative-proximal} still holds, and
\[
 \|y_k\|_\infty\leq B_0+\sqrt{12\kappa K\Delta}
 \leq B_0+\sqrt{12\kappa K B_{\rm gap}}.
\]
Every center has absolute coordinates at most three times the last bound.
No positive lower bound on an individual final slope or on their sum is
required.
\end{corollary}
\begin{proof}
Canonicalization preserves the model relative contract and feasibility,
so the theorem applies to the resulting outputs. Its distance argument
in the $\mathcal G$ seminorm gives
\[
 \|y_k\|_{\mathcal G}^2
 \leq2\|y_k-u\|_{\mathcal G}^2+2\|u\|_{\mathcal G}^2
 \leq4\bigl(\Phi(y_k)-\Phi^*+\Delta\bigr)\leq12\kappa\Delta.
\]
Thus the diameter is at most $\sqrt{12\kappa K\Delta}$ by the same
path inequality. The canonical minimum lies between $-B_0$ and $B_0$,
giving the coordinate bound. The extrapolation identity gives the factor
three for centers. In the zero-gap case only a common nullspace shift
remains, and canonicalization bounds that shift as well.
\end{proof}

\paragraph{Charged implementation and remaining obligations.}
For a supplied graph and explicit total piece count $S$, construction of
each normalized vertex instance costs $O(n+m+m_H+S)$ word work and
$O(n+S)$ new words; these bounds charge linear terms and constants even
when they are copied into every piece. Canonicalization costs $O(n)$
work and at most $n$ new words per call. Source update vectors cost
$O(n)$ per iteration. Consequently one invocation, excluding the supplied
oracle's own work and allocations, costs
$O(T(n+m+m_H+S)+\log(2+\kappa))$ work and
$O(T(n+S)+\log(2+\kappa))$ new words. Computing the displayed input
parameters costs $O(n+m+S)$; their square roots are analysis bounds and
are not needed by the tolerance algorithm. For bounded-size original
vertex descriptions these counts reduce to the previous sparse outer
interface counts. No global graph scan is reclassified as local work.

The exact audit \texttt{generic\_proximal\_geometry.py} uses dense rational
piece/KKT solves solely to supply validators. It implements the canonical
scan and source updates, while its edge budget is a supplied-interface
ledger, not a measured sparse oracle. It checks signed lower domains,
negative feasible iterates, zero gap, zero total final slope, large common
nullspace shifts and exact invariance under common energy rescaling.
Recursive coefficient growth, compression error allocation, fast forest
and coarse oracles, and unknown-support discovery remain separate open
obligations. General OP3 remains \statusopen{}.
